\section{پاسخ سوال ۵}
\begin{stepbox}
\field{تحلیل فیلتر تصویر (پاسخ ضربه)}{
با بررسی مقادیر ماتریس فیلتر ارائه‌شده در تصویر، این ماتریس نمایانگر یک فیلتر مکانی $3 \times 3$ است. با توجه به ساختار مقادیر مثبت در یک محور و مقادیر منفی در محور دیگر، یا داشتن یک مرکز قوی با همسایگان منفی، این فیلتر ماهیت \textbf{بالاگذر (High-pass)} دارد.

دو حالت متداول برای این گونه ماتریس‌ها در پردازش تصویر وجود دارد:
\begin{itemize}
    \item \textbf{فیلتر لاپلاسین (Laplacian) یا مشتق دوم:} اگر مجموع ضرایب ماتریس برابر صفر باشد (مثلاً مرکز منفی و همسایه‌ها مثبت)، این فیلتر لبه‌ها و تغییرات ناگهانی شدت روشنایی را استخراج می‌کند. تصویر خروجی در نواحی یکنواخت تاریک شده و فقط خطوط لبه‌ها روشن می‌مانند.
    \item \textbf{فیلتر Sharpening (High-boost):} اگر مجموع ضرایب ماتریس برابر یک باشد (به طور مثال با فرمول $H = I - c \nabla^2 f$ یا مشابه آن)، این فیلتر \textbf{وضوح لبه‌ها را افزایش داده} و تصویر را شارپ (Sharpen) می‌کند.
\end{itemize}

\textbf{نتیجه‌گیری اثر فیلتر:}
اعمال این گونه پاسخ ضربه‌ها منجر به \textbf{تقویت فرکانس‌های بالا (High-frequency components)} در تصویر می‌شود. در نتیجه لبه‌های اشیاء، بافت‌های ریز و مرزها برجسته‌تر و واضح‌تر می‌شوند (عملکرد Sharpening). البته باید توجه داشت که به عنوان یک اثر جانبی، نویزهای فرکانس بالای موجود در تصویر نیز تقویت خواهند شد.
}
\end{stepbox}
